% Part: first-order-logic
% Chapter: natural-deduction
% Section: derivations

\documentclass[../../../include/open-logic-section]{subfiles}

\begin{document}

\iftag{FOL}
      {\olfileid{fol}{ntd}{der}}
      {\olfileid{pl}{ntd}{der}}

\olsection{\usetoken{P}{derivation}}

\begin{explain}
موږ وويل چې فرضيه څه ده او د استنتاج قاعدې مو ورکړې۔ په فطري
استنتاج کښې !!^{derivation}s له دغو څخه په استقرايي ډول جوړېږي: هر
!!{derivation} يا يوازې يوه فرضيه وي، يا له يو، دوو يا دريو
!!{derivation}s او ورپسې له يوه سم استنتاج څخه جوړ وي۔
\end{explain}

\begin{defn}[!!^{derivation}]
\Article{derivation} له فرضيو~$\Gamma$ څخه د !!a{sentence}~$!A$ يو \emph{!!{derivation}}
د !!{sentence}s يوه متناهي ونه ده چې لاندې شرطونه پوره کوي:
\begin{enumerate}
\item د ونې تر ټولو پورتنۍ !!{sentence}s يا په~$\Gamma$ کښې وي يا
  په ونه کښې د يوه استنتاج په وسيله !!{discharged} شوې وي۔
\item د ونې تر ټولو لاندې !!{sentence}~$!A$ ده۔
\item په ونه کښې هره !!{sentence}، پرته له لاندېنۍ جملې~$!A$
  څخه، د استنتاج د يوې قاعدې د سمې کارونې مقدمه وي چې پايله ئې په
  ونه کښې نېغ د هغې !!{sentence} لاندې ولاړه وي۔
\end{enumerate}
بيا وايو چې $!A$ د !!{derivation} \emph{پايله} او~$\Gamma$ ئې
!!{undischarged} فرضيې دي۔

که له~$\Gamma$ څخه د~$!A$ کوم !!a{derivation} وي، وايو چې $!A$
له~$\Gamma$ څخه \emph{!!{derivable}} دے، يا په نښو کښې:
$\Gamma \Proves !A$۔ که د~$!A$ داسې !!a{derivation} وي چې هره
فرضيه پکې !!{discharged} شوې وي، نو~$\Proves !A$ ليکو۔
\end{defn}

\begin{ex}
هره فرضيه په يوازې ځان !!a{derivation} ده۔ نو مثلاً $!A$ په يوازې
ځان !!a{derivation} دے او $!B$ هم په يوازې ځان همداسې دے۔
له دغو څخه د، مثلاً، $\Intro{\land}$ قاعدې په کارولو يو نوے
!!{derivation} اخستلے شو:
\begin{prooftree}
\AxiomC{$!A$}
\AxiomC{$!B$}
\RightLabel{\Intro{\land}}
\BinaryInfC{$!A \land !B$}
\end{prooftree}
دا قاعدې عامې دي: پکې $!A$ او~$!B$ په هر ډول !!{sentence}s، مثلاً
په~$!C$ او~$!D$، بدلولے شو۔ بيا به پايله $!C \land !D$ وي؛ نو
\begin{prooftree}
\AxiomC{$!C$}
\AxiomC{$!D$}
\RightLabel{\Intro{\land}}
\BinaryInfC{$!C \land !D$}
\end{prooftree}
يو سم !!{derivation} دے۔ البته فرضيې سره اړولے هم شو، څو $!D$ د~$!A$
او $!C$ د~$!B$ ونډه واخلي۔ له دې امله
\begin{prooftree}
\AxiomC{$!D$}
\AxiomC{$!C$}
\RightLabel{\Intro{\land}}
\BinaryInfC{$!D \land !C$}
\end{prooftree}
هم يو سم !!{derivation} دے۔

اوس بله قاعده، مثلاً $\Intro{\lif}$، کارولے شو۔ دا اجازه راکوي چې
يو شرطي قضيه پايله کړو او هره هغه فرضيه !!{discharge} کړو چې له
دغه شرطي قضيې له مقدم سره يو شان وي۔ نو لاندې دواړه سم
!!{derivation}s دي:
\begin{prooftree}
\AxiomC{$\Discharge{!C}{1}$}
\AxiomC{$!D$}
\RightLabel{\Intro{\land}}
\BinaryInfC{$!C \land !D$}
\DischargeRule{\Intro{\lif}}{1}
\UnaryInfC{$!C \lif (!C \land !D)$}
\DisplayProof\bottomAlignProof
\AxiomC{$!C$}
\AxiomC{$\Discharge{!D}{1}$}
\RightLabel{\Intro{\land}}
\BinaryInfC{$!C \land !D$}
\DischargeRule{\Intro{\lif}}{1}
\UnaryInfC{$!D \lif (!C \land !D)$}
\end{prooftree}
دوئ په ترتيب سره ښيي چې $!D \Proves !C \lif (!C \land !D)$ او
$!C \Proves !D \lif (!C \land !D)$۔

راياد کړئ چې د فرضيو ايستل اجازه ده، شرط نۀ دے: لازم نۀ ده چې
فرضیې وباسو۔ په ځانګړي ډول، يوه قاعده هله هم کارولے شو چې فرضيې په
!!{derivation} کښې موجودې نۀ وي۔ مثلاً لاندې اشتقاق روا دے، که څه
هم د !!{discharged} کېدو دپاره هېڅ فرضيه~$!A$ نۀ شته:
\begin{prooftree}
  \AxiomC{$!B$}
  \DischargeRule{\Intro{\lif}}{1}
  \UnaryInfC{$!A \lif !B$}
\end{prooftree}
\end{ex}

\end{document}
